\documentclass[11pt]{article}

\usepackage[a4paper,margin=30mm]{geometry}
\usepackage{amsmath,amssymb,amsthm}
\usepackage[hidelinks]{hyperref}
\usepackage{microtype}

\newtheorem{theorem}{Theorem}
\newtheorem{lemma}[theorem]{Lemma}
\newtheorem{proposition}[theorem]{Proposition}
\theoremstyle{definition}
\newtheorem{definition}[theorem]{Definition}
\newtheorem{remark}[theorem]{Remark}

\newcommand{\Fp}{\mathbb F_p}
\newcommand{\Zp}{\mathbb Z_p}
\newcommand{\alphaG}{\alpha}

\title{From interval palettes to dense nonsumsets in prime fields}
\author{Project proof artifact for AIM Problem 2.7}
\date{26 August 2026}

\begin{document}
\maketitle

\begin{abstract}
For a prime $p$, let $M(p)$ be the greatest cardinality of a subset of
$\mathbb F_p$ which is not $B+B$ for any $B\subseteq\mathbb F_p$, with equal
summands allowed.  We prove
\[
  M(p)=p-p^{1/2+o(1)}.
\]
The proof combines Noga Alon's universal large-set theorem with the
interval-palette theorem in Shouqiao Wang's pinned proposed solution of
Erd\H{o}s Problem 788.  We transfer two interval palettes to consecutive
charts of $\mathbb Z_p$, bound every possible self-sumset root by an
independence number, and use a strict finite count to choose an unrealized
output.  All imported statements and new reductions are identified
explicitly.  Wang's theorem is machine-checked but recent and not
peer-reviewed, so the result is presented as a candidate proof for expert
review rather than as a community-accepted resolution.
\end{abstract}

\section{Introduction}

The problem, recorded as AIM Problem 2.7 by B. Green \cite{AIM2004}, asks for
the largest size of a subset of a prime field which is not an unrestricted
self-sumset.  Theorem~\ref{thm:main} determines the exponent of the missing
part: the extremal nonsumsets have $p^{1/2+o(1)}$ missing elements.  The
argument is short once one has a suitable low-score graph on the additive
group, but the passage from the available interval graph to the cyclic group
and the requirement of an \emph{exact} sumset output both need justification.

\paragraph{Contribution and dependency boundary.}
There are two imported results.  Alon's published theorem gives the lower
bound on the number of missing elements.  Wang's interval-palette theorem
gives graphs whose palette size plus independence number is
$N^{1/2+o(1)}$.  The new part of this paper consists of the normalization of
Wang's theorem, the two-chart cyclic transfer, and the exact-output deletion
argument.  The mathematical deduction is complete relative to the two stated
inputs.  Appendix~\ref{sec:evidence} records separately why the recent Wang
input is treated as machine-checked evidence rather than as an already
accepted literature result.

\paragraph{Proof roadmap.}
The proof has five steps.
\begin{enumerate}
  \item Alon's theorem shows that deleting at most
  $\frac1{4000}\sqrt{p/\log p}$ elements always leaves a self-sumset.
  \item Wang's theorem supplies an interval graph whose palette size $e$ and
  independence number $r$ have sum $p^{1/2+o(1)}$ after rescaling.
  \item Two consecutive interval charts cover $\mathbb Z_p$ and transfer
  these palettes to a cyclic graph $\Gamma_E$ without increasing the score
  $H=e+r$ beyond $p^{1/2+o(1)}$.
  \item If $(B+B)\cap E=\varnothing$, then the distinct elements of $B$ form
  an independent set in $\Gamma_E$, so every possible root has $|B|\le r$.
  \item There are more $4H$-element deletion sets outside $E$ than there are
  roots of size at most $r$.  One deletion therefore produces an exact output
  which is not $B+B$ for any $B$.
\end{enumerate}
The lower and upper estimates for the number of missing elements then have
the same exponent.

\section{Statement and notation}

For $B\subseteq\Fp$, define the unrestricted self-sumset
\[
  B+B=\{b_1+b_2:b_1,b_2\in B\}.
\]
The empty root is allowed, and $b_1=b_2$ is allowed.  Define
\[
  M(p)=\max\{|A|:A\subseteq\Fp\text{ and }A\ne B+B
                    \text{ for every }B\subseteq\Fp\}
\]
and put $d(p)=p-M(p)$.  Thus $d(p)$ is the strict complement threshold,
not the inclusive threshold used in some earlier papers.

\begin{theorem}[AIM Problem 2.7]\label{thm:main}
As $p$ tends to infinity through the primes,
\[
  d(p)=p^{1/2+o(1)}
  \qquad\text{and}\qquad
  M(p)=p-p^{1/2+o(1)}.
\]
\end{theorem}

All logarithms are natural.  We identify $\Fp$ with the additive group
$\Zp$ and with its standard representatives $\{0,\ldots,p-1\}$.

\section{Two imported results}

The lower estimate follows from the following exact inclusive statement.

\begin{theorem}[Alon, Theorem 2.1 \cite{Alon2007}]\label{thm:alon}
For every sufficiently large prime $p$, if $F\subseteq\Zp$ and
\[
  |F|\le \frac{1}{4000}\sqrt{\frac p{\log p}},
\]
then $\Zp\setminus F=A+A$ for some $A\subseteq\Zp$.
\end{theorem}

We next state exactly the consequence of Wang's interval theorem which is
used here.  For $N\ge2$ and
$Q\subseteq\{1,\ldots,2N-3\}$, let $H_Q$ be the simple graph on
\[
  [N]_0=\{0,1,\ldots,N-1\}
\]
in which distinct $x,y$ are adjacent precisely when $x+y\in Q$.

\begin{theorem}[Interval-palette theorem \cite{Wang2026}]\label{thm:interval}
There is an absolute constant $C>0$ such that, for every sufficiently large
integer $N$, there is $Q_N\subseteq\{1,\ldots,2N-3\}$ satisfying
\[
  |Q_N|+\alphaG(H_{Q_N})
  \le (N+1)^{1/2+\Delta_N},
  \qquad
  \Delta_N=C\left(\frac{\log\log(N+1)}{\log(N+1)}\right)^{1/3}.
\]
\end{theorem}

For completeness, this is the exact normalization of the cited result.
Wang writes
\[
  I_n=(n,2n)\cap\mathbb N=\{n+1,\ldots,2n-1\}
\]
and proves both the min--max identity
\[
  f(n)=\min_{D\subseteq J_n}\bigl(|D|+\alpha(G_D)\bigr),
\]
with the minimum unchanged when $D$ is restricted to attainable sums of two
distinct elements of $I_n$, and the upper bound
\[
  f(n)\le n^{1/2+C(\log\log n/\log n)^{1/3}}.
\]
Set $n=N+1$, translate $I_n$ by $-(n+1)$, and translate every palette sum by
$-(2n+2)$.  The vertex set becomes $[N]_0$, the attainable distinct sums
become exactly $\{1,\ldots,2N-3\}$, and the graph and its independence number
are preserved.  This gives Theorem~\ref{thm:interval} without an asymptotic
change of parameters.

\section{Transfer to a cyclic palette}

For $E\subseteq\Zp$, let $\Gamma_E$ be the simple graph on $\Zp$ in which
distinct $x,y$ are adjacent exactly when $x+y\in E$.

\begin{lemma}[Root-size reduction]\label{lem:root-bound}
For any $E,B\subseteq\Zp$, if $(B+B)\cap E=\varnothing$, then $B$ is
independent in $\Gamma_E$.  In particular,
\[
  |B|\le\alphaG(\Gamma_E).
\]
\end{lemma}

\begin{proof}
If distinct $x,y\in B$ were adjacent, then
$x+y\in(B+B)\cap E$, contrary to the hypothesis.  Thus no two distinct
vertices of $B$ are adjacent.  Notice that equal summands remain in $B+B$;
the conclusion only needs the distinct-pair edges of the simple graph.
\end{proof}

\begin{lemma}[Two-chart transfer]\label{lem:transfer}
For every sufficiently large odd prime $p$, there is $E\subseteq\Zp$ such
that, with
\[
  e=|E|,\qquad r=\alphaG(\Gamma_E),\qquad H=e+r,
\]
one has
\[
  2\sqrt p-1\le H\le p^{1/2+\eta_p}
\]
for a sequence $\eta_p\to0$.
\end{lemma}

\begin{proof}
Put
\[
 N_0=\frac{p+1}{2},\quad N_1=\frac{p-1}{2},\quad
 a_0=0,\quad a_1=N_0,
\]
and partition $\Zp$ into the two standard intervals
\[
 V_i=a_i+\{0,\ldots,N_i-1\}\qquad(i=0,1).
\]
Choose $Q_{N_i}$ from Theorem~\ref{thm:interval} and define
\[
 E_i=\{2a_i+q\pmod p:q\in Q_{N_i}\},
 \qquad E=E_0\cup E_1.
\]
Because $Q_{N_0}\subseteq\{1,\ldots,p-2\}$ and
$Q_{N_1}\subseteq\{1,\ldots,p-4\}$, reduction modulo $p$ is injective on
each palette.  Cross-palette collisions are allowed and only decrease $|E|$.
Thus
\[
 e\le |Q_{N_0}|+|Q_{N_1}|.                 \label{eq:palette-size}
\]

Let $R$ be independent in $\Gamma_E$.  For fixed $i$, define
\[
 R_i=\{u\in[N_i]_0:a_i+u\in R\}.
\]
If distinct $u,v\in R_i$ satisfied $u+v\in Q_{N_i}$, then
\[
 (a_i+u)+(a_i+v)\equiv2a_i+u+v\pmod p\in E_i\subseteq E,
\]
contradicting independence.  Hence $R_i$ is independent in $H_{Q_{N_i}}$.
The two charts partition $\Zp$, so
\[
 r\le\alphaG(H_{Q_{N_0}})+\alphaG(H_{Q_{N_1}}).
                                                        \label{eq:independence}
\]
Let $\varepsilon_p=\max_i\Delta_{N_i}$ and
$\eta_p=\varepsilon_p+\log 2/\log p$.
Equations~\eqref{eq:palette-size} and \eqref{eq:independence}, together with
Theorem~\ref{thm:interval}, give
\[
 H\le\sum_{i=0}^1(N_i+1)^{1/2+\Delta_{N_i}}
   \le2p^{1/2+\varepsilon_p}=p^{1/2+\eta_p}.
\]
Since $N_i\to\infty$, we have $\eta_p\to0$.

For the lower estimate, each vertex of $\Gamma_E$ has degree at most $e$:
for fixed $x$ and fixed $s\in E$, at most one neighbor can satisfy $x+y=s$.
A graph of maximum degree $e$ has an independent set of size at least
$p/(e+1)$ by greedy coloring.  Therefore
\[
 H=e+r\ge e+\frac p{e+1}
   =(e+1)+\frac p{e+1}-1\ge2\sqrt p-1.
\]
This argument uses only distinct-pair edges, so no diagonal convention has
been changed.
\end{proof}

\section{The exact-output deletion}

\begin{lemma}[Strict finite counting gap]\label{lem:count}
With $e,r,H,\eta_p$ as in Lemma~\ref{lem:transfer}, for every sufficiently
large $p$,
\[
  3\le H,\qquad 4H\le p-e,
\]
and
\[
  \binom{p-e}{4H}>\sum_{j=0}^{r}\binom pj.          \label{eq:counting-gap}
\]
\end{lemma}

\begin{proof}
Take $p$ large enough that
\[
 \eta_p\le\frac1{16},\qquad
 \frac{\log8}{\log p}\le\frac1{16},\qquad
 5p^{9/16}\le p.                                  \label{eq:large-p}
\]
Lemma~\ref{lem:transfer} gives
\[
 2\sqrt p-1\le H\le p^{1/2+\eta_p}\le p^{9/16}.
\]
Hence $H\ge3$ eventually and, because $e\le H$,
$e+4H\le5p^{9/16}\le p$.

For integers $0<k\le n$,
\[
 \binom nk=\prod_{i=0}^{k-1}\frac{n-i}{k-i}
 \ge\left(\frac nk\right)^k,
\]
because $(n-i)/(k-i)\ge n/k$ is equivalent to $i(n-k)\ge0$.
Apply this with $n=p-e$ and $k=4H$.  Since $e\le H\le p^{9/16}\le p/2$
eventually,
\[
 \log\binom{p-e}{4H}
 \ge4H\log\frac{p-e}{4H}
 \ge4H\log\frac p{8H}.
\]
Using \eqref{eq:large-p} and $H\le p^{9/16}$, the last logarithm is at least
$\frac38\log p$, and therefore
\[
 \binom{p-e}{4H}\ge p^{3H/2}.                    \label{eq:binomial-lower}
\]
On the other hand,
\[
 \sum_{j=0}^{r}\binom pj
 \le\sum_{j=0}^{r}p^j
 \le p^{r+1}\le p^{H+1}\le p^{4H/3}<p^{3H/2},
\]
where $H+1\le4H/3$ follows from $H\ge3$.  Combining this strict inequality
with \eqref{eq:binomial-lower} proves \eqref{eq:counting-gap}.
\end{proof}

\begin{lemma}[Unrealized output]\label{lem:deletion}
For every sufficiently large prime $p$, there is $T\subseteq\Zp\setminus E$
with $|T|=4H$ such that
\[
 A_p=\Zp\setminus(E\cup T)
\]
is not $B+B$ for any $B\subseteq\Zp$.  Moreover,
\[
 |\Zp\setminus A_p|\le5p^{1/2+\eta_p}.
\]
\end{lemma}

\begin{proof}
Let $\mathcal T$ be the family of all $4H$-element subsets of
$\Zp\setminus E$.  By Lemma~\ref{lem:count},
\[
 |\mathcal T|=\binom{p-e}{4H}>
 \sum_{j=0}^{r}\binom pj.                         \label{eq:deletion-count}
\]
Call $T\in\mathcal T$ realized if
\[
 B+B=\Zp\setminus(E\cup T).                       \label{eq:realized-output}
\]
for some $B\subseteq\Zp$.  Equation~\eqref{eq:realized-output} implies
$(B+B)\cap E=\varnothing$.  Lemma~\ref{lem:root-bound} therefore gives
$|B|\le r$.  There are at most the right-hand side of
\eqref{eq:deletion-count} such roots.  In addition, a fixed root determines
at most one deletion set, because \eqref{eq:realized-output} forces
\[
 T=(\Zp\setminus(B+B))\setminus E.                \label{eq:forced-deletion}
\]
Consequently \eqref{eq:deletion-count} leaves some $T$ unrealized.  The
corresponding $A_p$ is a nonsumset.  Finally, $E\cap T=\varnothing$ and
$e\le H$, so
\[
 |\Zp\setminus A_p|=e+4H\le5H\le5p^{1/2+\eta_p}.
\]
Equal summands cause no gap: they remain part of
\eqref{eq:realized-output}, while distinct pairs provide the necessary
independence condition and \eqref{eq:forced-deletion} uses the full sumset.
\end{proof}

\section{Proof of the main theorem}

\begin{proof}[Proof of Theorem~\ref{thm:main}]
Set
\[
 q_p=\left\lfloor\frac1{4000}\sqrt{\frac p{\log p}}\right\rfloor.
\]
If $S\subseteq\Zp$ has $|S|\ge p-q_p$, then its complement satisfies the
inclusive hypothesis of Theorem~\ref{thm:alon}; hence $S$ is a self-sumset.
No nonsumset has size at least $p-q_p$, and integrality gives
\[
 M(p)\le p-q_p-1,
 \qquad d(p)\ge q_p+1>\frac1{4000}\sqrt{\frac p{\log p}}.
                                                        \label{eq:lower-d}
\]

Lemma~\ref{lem:deletion} supplies a nonsumset $A_p$ with complement at most
$5p^{1/2+\eta_p}$.  Hence
\[
 d(p)=p-M(p)\le p-|A_p|\le5p^{1/2+\eta_p}.         \label{eq:upper-d}
\]
Taking logarithms in \eqref{eq:lower-d} and \eqref{eq:upper-d} yields
\[
 \frac12-\frac{\log\log p}{2\log p}-\frac{\log4000}{\log p}
 <\frac{\log d(p)}{\log p}
 \le\frac12+\eta_p+\frac{\log5}{\log p}.
\]
Both outer expressions tend to $1/2$.  Therefore
$d(p)=p^{1/2+o(1)}$, and the defining identity $M(p)=p-d(p)$ gives
$M(p)=p-p^{1/2+o(1)}$.
\end{proof}

\appendix

\section{Verification and evidence boundary}\label{sec:evidence}

The imported Wang repository was pinned at commit
\[
\texttt{d28713ac8245ca86a686b8c67370a8d19d81b242}.
\]
Its Lean 4.27 development used Mathlib commit
\[
\texttt{a3a10db0e9d66acbebf76c5e6a135066525ac900}
\]
and completed all 7,936 local build jobs.  The final theorem depends only on
\texttt{propext}, \texttt{Classical.choice}, and \texttt{Quot.sound}; the
audited source contains no \texttt{sorry}, \texttt{axiom}, or \texttt{admit}.
The interval normalization above was checked against both the paper and the
formal definitions.

The new cyclic transfer and deletion proof passed a separate QED xhigh
structural audit, a detailed step-by-step audit, deterministic finite
falsification checks, and final verdict \texttt{DONE}.  These checks establish
the repository's proof claim relative to the pinned imported theorem.  They do
not establish novelty, peer review, or community acceptance.  In particular,
the source itself is titled a proposed solution and the public Erd\H{o}s
problem page remained marked open at the audit date.

\begin{thebibliography}{9}

\bibitem{AIM2004}
E. Croot and V. F. Lev (collectors),
\emph{Problems Presented at the Workshop on Recent Trends in Additive
Combinatorics}, AIM workshop list, Problem 2.7 (B. Green), 2004.
\url{https://aimath.org/WWN/additivecomb/additivecomb.pdf}.

\bibitem{Alon2007}
N. Alon,
\emph{Large sets in finite fields are sumsets},
J. Number Theory 126 (2007), 110--118.
\url{https://web.math.princeton.edu/~nalon/PDFS/sumset.pdf}.

\bibitem{Wang2026}
S. Wang,
\emph{A Proposed Solution to Erd\H{o}s Problem 788},
pinned repository artifact, commit
\texttt{d28713ac8245ca86a686b8c67370a8d19d81b242}, 2026.
\url{https://github.com/ShouqiaoW/erdos/tree/d28713ac8245ca86a686b8c67370a8d19d81b242/788}.

\end{thebibliography}

\end{document}
